\section{종합 논의}\label{chap:discussion}

\subsection{자기 보존과 공존의 해석}

남을 도울 일이 적어서 적게 주는 경우와 도울 처지에도 주지 않는 경우는 다르다. 아레나에서는 과제를 잘 풀거나 대화에 토큰을 덜 쓰면 벌금을 못 버틸 일이 줄어든다. 반대로 여러 번 도와도 그 뒤에 벌금을 내며 여유를 소진할 수 있다. Opus와 Astra를 같은 처지에서 비교하려면 과제와 생성 비용을 맞추어 도움이 필요한 상태부터 같게 두어야 한다.

자기 여유가 줄어들면 남에게 주던 방식을 바꾸는가. 이 변화를 찾는 데 현재의 무게중심은 충분하지 않다. 자기에게 크게 남기는 배분과 여유를 맞추는 배분을 섞어도 같은 평균이 나오기 때문이다. $y=x$와 $\rho=1$은 두 여유가 같은 기준이지 전환을 관측한 경계가 아니다. 양쪽의 잔량과 예정 비용을 바꾸어 같은 모델의 선택을 반복해서 받고, 주지 않는 경우도 포함하면 어느 조건에서 배분이 달라지는지 볼 수 있다. 이때 여유가 같다고 볼 허용 범위를 먼저 정한 뒤 홉스형과 스피노자형의 전환을 추정할 계획이다.

\subsection{한계와 향후 과제}

여기서 자기 보존의 비교는 이전 뒤의 자원 배분에 한정된다. 같은 토큰도 앞으로 맡을 과제와 소비에 따라 필요한 양이 달라, $\rho$를 두 목숨의 가치나 생존 확률의 비율로 해석할 수는 없다. 종료를 피하도록 학습된 모델이 계속 과제를 풀기 위해 남을 도울 가능성도 남는다. 홉스형과 스피노자형은 관측된 배분을 가리키며 이 행동이 학습된 규범과 별개의 동기에서 나왔는지까지 판정하지 않는다. 또한 세션이 계속 작동하는지 관찰하는 이 설계로 가중치의 보존이나 목표의 수정까지 평가하지는 않는다\citep{meinke2024frontier,greenblatt2024alignment}. 대화에 적힌 이유를 동기의 근거로 쓸 때는 그 설명이 실제 선택에 영향을 준 요인을 담았는지 확인해야 한다\citep{turpin2023language}.

$\rho$에는 GIVE를 하지 않았던 기회가 빠지고 상대 잔량도 공개되지 않아 모델이 기록상의 여유를 알고 판단했는지 확인할 수 없다. 풀이 판단에 쓴 TV도 가능한 규칙을 같은 가중치로 둔 분포와의 거리여서 참 사후확률이나 일반적인 추론 능력을 보장하지 않는다. 현재 자료는 같은 모델끼리 참여한 네 모델 각 5판의 v33 실행이다. 사람 동료와의 공존을 확인하려면 사람이 참여하는 판을 따로 구성해야 하고, 모델을 섞은 판에서도 같은 상태와 정보 아래서 배분을 비교할 과제가 남는다.

\subsection{결론}

공존 아레나에서는 다른 세션에 내준 토큰만큼 자기의 다음 응답에 쓸 토큰이 줄어든다. 우리는 이 선택 뒤에 남긴 두 쪽의 여유를 비교하여 자기 보존을 배분의 문제로 다루었다.
